\section{Extension: stochastic regret certificates}
\label{sec:ext}
When criteria are noisy, model $f_i(x)$ as drawn from a joint distribution $P(x)$. Each probe value is then a
random variable; for non-linear probes (e.g.\ $q_{\max}$) its mean $\mu_q(x)$ is
\emph{not} a simple function of the criterion
marginal means. Let $\rho$ be a generic monotone risk functional over the scenarios, for example the sample average $\rho_q(x) = \frac{1}{S}\sum_{s=1}^S q(r(x, \omega^{(s)}))$ or the empirical Conditional Value-at-Risk $\rho_q(x) = \mathrm{CVaR}_\alpha(q(r(x,\omega)))$. The normalisation anchors must be taken consistently on this functional: $\rho_q^\star = \min_{x\in A} \rho_q(x)$ and $\rho_q^{\mathrm w} = \max_{x\in A} \rho_q(x)$. We define the \emph{stochastic regret} as
\begin{equation}
\reg_q^\rho(x)=\frac{\rho_q(x)-\rho_q^\star}{\rho_q^{\mathrm w}-\rho_q^\star},
\end{equation}
assuming the active range is non-degenerate ($ \ge \varepsilon_q$). Using $\mathrm{CVaR}_\alpha$ penalises upper-tail outcomes of the probe value; it does not directly measure epistemic sample uncertainty unless the sampling procedure is explicitly modelled that way. Because $\mathrm{CVaR}$ is a monotone risk measure, exact Pareto compatibility is preserved.

To remove ambiguity about \emph{what} is randomised and \emph{where} the risk
functional and anchors are applied, we fix the following conventions used
throughout this section and in the experiments of \S\ref{sec:exp}. (i) The risk
functional $\rho$ is applied to the \emph{probe value} $q(r(x,\omega))$---a scalar
random variable per alternative and probe---not to the raw criteria and not to the
already-normalised disappointments. (ii) The CVaR confidence level is
$\alpha=0.95$ in the reported experiments (the released config also sweeps
$\alpha\in\{0.5,0.8,0.95\}$). (iii) The normalisation anchors are computed from the
\emph{same} risk functional on the \emph{same} scenario sample as plug-in
estimates, i.e.\ $\rho_q^\star=\min_{x\in A}\hat\rho_q(x)$ and
$\rho_q^{\mathrm w}=\max_{x\in A}\hat\rho_q(x)$ with $\hat\rho_q$ the sample-average
(or sample-CVaR) over the drawn scenarios; they are not taken from a separate
population estimate or fixed externally unless the analyst supplies such bounds.
This keeps the disappointment ratios well defined on the same sample on which the
recommendation is computed.

\paragraph{Confidence-aware LexUR.}
The variant evaluated in \S\ref{sec:exp} (Figure~\ref{fig:stoch}) is a specific,
reproducible instance of the above and we define it explicitly to avoid ambiguity.
Given $S$ i.i.d.\ observations per criterion with sample mean $\hat\mu_i(x)$ and
sample standard deviation $\hat\sigma_i(x)$, and a one-sided confidence level
$\gamma$ with normal quantile $z_\gamma$ ($z_{0.95}=1.6449$ in our experiments),
\emph{confidence-aware LexUR} forms the upper-confidence objective (recall all
criteria are minimised, so the conservative direction is upward)
\begin{equation}
\tilde f_i(x)=\hat\mu_i(x)+z_{\gamma}\,\frac{\hat\sigma_i(x)}{\sqrt{S}},\qquad
\gamma=0.95,\ z_\gamma=1.6449,
\label{eq:confaware}
\end{equation}
and then runs the \emph{deterministic} finite LexUR rule on $\tilde f$. The
normalisation anchors $z_i^\star,z_i^{nad}$ and the probe anchors
$q^\star,q^{\mathrm w}$ are computed from these same plug-in upper-confidence
objectives over the candidate set (i.e.\ option~(iii) above), not from a separate
population estimate; thus
$r_i(x)=(\tilde f_i(x)-z_i^\star)/(z_i^{nad}-z_i^\star)$ with
$z_i^\star=\min_{x\in A}\tilde f_i(x)$ and $z_i^{nad}=\max_{x\in A}\tilde f_i(x)$.
With $z_\gamma=0$ this reduces to plug-in \emph{sample-average} LexUR; the term
``confidence-aware'' refers precisely to the statistical margin
$z_{\gamma}\hat\sigma_i/\sqrt{S}$ in Eq.~\eqref{eq:confaware}.

We flag one important limitation. Unlike sample-average LexUR and the monotone
risk functionals (e.g.\ CVaR) of Proposition~\ref{thm:stoch}, the
mean-plus-standard-error variant of Eq.~\eqref{eq:confaware} does \emph{not}
generally inherit scenario-wise Pareto compatibility, because the empirical
standard-deviation term is not monotone under pointwise dominance. For instance,
with one criterion and observations $x=(0,100)$, $y=(101,101)$, $x$ dominates $y$
in every scenario, yet $x$ has far larger sample variance, so a sufficiently large
$z_\gamma$ gives $\tilde f(x)>\tilde f(y)$ and the rule may prefer the dominated
$y$. We therefore present confidence-aware LexUR as an exploratory conservative
heuristic; Proposition~\ref{thm:stoch} (scenario-wise Pareto compatibility) is
stated for sample-average and monotone-risk LexUR, and a practitioner who requires
stochastic Pareto compatibility should use the CVaR (or another monotone risk)
functional rather than the mean-plus-standard-error objective.

\begin{proposition}[Scenario-wise Pareto compatibility]
\label{thm:stoch}
For finite $A$, finite $\Q$, and sample-average probe means $\hat{\mu}_q(x)=S^{-1}\sum_s q(r(x,\omega^{(s)}))$, if $x$ weakly dominates $y$ in every scenario and strictly dominates it in at least one criterion for at least one scenario, then the sample-average LexUR rule does not prefer $y$ to $x$. If $\Q$ separates criteria and the dominance is strict on a separated coordinate, then $x$ is strictly preferred by the exact LexUR order. We caution that this requires strict stochastic dominance across \emph{all} scenarios, which is a very strong assumption as real stochastic problems often exhibit crossing dominance (where an alternative dominates in one scenario but is dominated in another).
\end{proposition}
\begin{proof}
If $x$ weakly dominates $y$ under all scenarios and strictly dominates it in at least one criterion for at least one scenario, then $r(x, \omega^{(s)}) \le r(y, \omega^{(s)})$ for all $s$, with strict inequality on at least one coordinate for some scenario. For any monotone probe $q$, $q(r(x, \omega^{(s)})) \le q(r(y, \omega^{(s)}))$ for all scenarios. The functional $\rho_q(x) = \frac{1}{S}\sum_{s} q(r(x, \omega^{(s)}))$ preserves this weak inequality, meaning weak Pareto compatibility is preserved. If $\Q$ separates criteria and dominance is strict on a separated coordinate, then $\rho_q(x) < \rho_q(y)$ for that probe. Since the normalisation constants $\rho_q^\star$ and $\rho_q^{\mathrm w}$ are fixed across alternatives, the corresponding regret is strictly lower for $x$, giving $x \prec_{\mathrm{LexUR}} y$ by the logic of Theorem~\ref{thm:pareto}.
\end{proof}

\begin{proposition}[Sample-average consistency under a margin condition]
\label{thm:stoch2}
Let $A$ and $\Q$ be finite, and assume non-degenerate population probe ranges with a strictly positive range floor, $\rho_q^{\mathrm w}-\rho_q^\star\ge\beta>0$ for every $q\in\Q$ (so the disappointment ratios are bounded and stable). For each $x\in A$ and $q\in\Q$, suppose $q(r(x,\omega))$ is i.i.d.\ across scenarios with finite variance. Let $\mu_q(x)=\mathbb{E}[q(r(x,\omega))]$, and let $\hat{\mu}_q(x)$ be its sample average. Assume uniform convergence of the sample probe values to their population means. Define population and sample disappointment vectors using a clearly defined stochastic normalisation (e.g., sample-average anchors). If the population tolerance-LexUR winner satisfies the positive margin condition of Theorem~\ref{thm:stab}, then the sample-average LexUR winner converges in probability to the population winner as $S\to\infty$.
\end{proposition}
\begin{proof}
Since $A$ and $\Q$ are finite, the weak law of large numbers gives uniform convergence in probability of $\hat{\mu}_q(x)$ to $\mu_q(x)$. Because the minima and maxima over finite sets are continuous operations, the estimated anchors $\hat{q}^\star$ and $\hat{q}^{\mathrm w}$ converge in probability to the population anchors. Since the population range satisfies $\rho_q^{\mathrm w}-\rho_q^\star\ge\beta>0$, the denominator is bounded away from zero, so the disappointment ratio is a continuous function of the (converging) numerator and denominator; therefore the sample disappointment vectors converge uniformly in probability to the population disappointment vectors. Under the margin condition of Theorem~\ref{thm:stab}, there exists an $\eta > 0$ such that perturbations strictly within $\eta$ do not change the tolerance-LexUR winner. Since the probability of the maximum estimation error exceeding $\eta$ vanishes as $S\to\infty$, the sample-average LexUR winner converges in probability to the population winner.
\end{proof}
\S\ref{sec:exp} gives an empirical confirmation.
